\chapter{总结与展望}\label{chap:introduction}

前面的章节详细介绍了开展本论文的背景与意义、国内外研究现状、面向网络流量的重复音视频文件实时检测系统的两个关键技术以及整个系统的设计和实现，本章对以上研究内容进行一次全面的总结，对研究过程中没有考虑过的问题进行详细分析，对下一步研究工作做出合理的展望。

\section{全文总结}
越来越多的人通过网络观看和下载音视频节目，这也使得音视频流量成为网络流量中占比最高的流量。随着音视频流量快速增长，对音视频流量的重复性检测也越来越受到广泛关注。在学术界，有大量的工作研究网络冗余发现，用以减少整个网络环境的资源开销。

面向数据链路层的网络冗余发现系统，面对的挑战是数据链路层网络的不稳定性和IP数据的比例下降的特性。链路层的冗余发现系统，主要部署在蜂窝数据网络和无线局域网。代表性的工作有\cite{Lumezanu2010The}和\cite{Halepovic2012}。

面向网络层的网络冗余发现系统，面对的挑战是检测的准确性和高效性。网络层的冗余发现系统通常是基于端到端的模型，通过不同的网络冗余发现算法对流量的冗余进行检测，其中较为出名的算法是PACK、MODP和拉宾指纹。代表性的工作有\cite{Spring2000}和\cite{Yu}。

面向应用层的网络冗余发现系统，面对的挑战是加密流量的兴起，其主要应对不同的应用场景如移动设备、企业内部网络、云计算等。代表性的工作有\cite{Aggarwal}和\cite{Feng2013How}。

本文的研究目标是针对网络中间件消耗大量资源处理重复音视频文件的问题，设计并实现面向网络流量的重复音视频实时检测系统。该系统有两个要求：一是检测的实时性。二是算法的准确性。

本文的两个工作基础，分别是流式模糊哈希和基于流式模糊哈希的大规模相似性查找系统，流式模糊哈希主要是在真实环境中生成音视频流量的摘要，作为基准用以生成文件编号。而相似性查找系统主要是为了计算两个流式模糊哈希值之间的编辑距离，评价两个文件之间的相似性。

整个重复音视频实时检测系统主要基于的两个技术：

1）预测性文件标识生成方法。通过重复音视频文件预测实验，传统的机器学习方法准确率和召回率结果不理想，不适用重复音视频文件实时检测。预测性文件标识生成方法采用信息熵和互信息作为评价标准，对生成文件标识的特征进行筛选，并在流量建立连接时，生成一个具有预测性的文件标识。 

2）基于流式模糊哈希的重复音视频检测方法。若文件标识可在数据库中查询到，则认为该音视频文件是重复的，同时通过流式模糊哈希值生成文件编号，通过两个指标验证文件标识的效果，分别为误报率漏报率。误报率漏报率的计算结果可用于指导文件标识的特征选择。通过实验得出了最优的文件标识候选特征，分别为Etag、MediaLen、Last-Modified、头部文件32K（其中Etag和Last-Modified任选其一）。

基于以上的两个技术，本文实现了重复音视频实时检测系统，该系统去重效果显著，文件数去重率约为60\%以上，带宽去重率约为20\%。经过该系统的流量内容分析模块特定音视频文件召回率高于95\%，未影响内容分析模块准确性。

本文研究的完成以下创新：

1）提出了一种预测性文件标识生成方法，该方法选取音视频流量头部的一些特征，以及文件头部的一些信息，通过熵和互信息对这些特征进行筛选，最后通过哈希算法生成一个具有预测性的文件标识。

2）提出了一种基于流式模糊哈希的重复音视频检测方法，在通过文件标识检测音视频文件重复性的同时，也会计算文件的流式模糊哈希值，用以生成文件编号，指导文件标识的生成。

3）实现了重复音视频实时检测系统，真实网络环境实验结果表明，该系统的去重机制未影响内容分析模块准确性，且去重效果显著。

\section{研究展望}

本文的研究工作存在以下几个方面需要进一步完善：

1）目前系统对于不同的文件标识，设计的生命周期为24小时。但测试中发现许多热点音视频的文件标识的生命周期可以延长到十几天甚至几个月，是否可以针对不同的文件，生成不同的标识生命周期？采用什么方法对文件标识的生命周期进行预测？

2）目前系统的文件标识评价模块，采用半离线的方式，抽取10\%的音视频流量，并计算其流式模糊哈希值，从而指导文件标号的生成。该模块应对概念漂移能力不理想，是否可以通过某些方法，实时地预测文件标识的准确性？

3）基于流式模糊哈希的大规模相似性查找系统在面对海量的流式模糊哈希查找时，耗费的内存高，无法在内存集约的环境下使用。是否可以对该查找系统进行改进，降低其内存消耗。